% Exact test of the compensated delta-method weight.  Work record only.
% No impossibility theorem about Sun's conjecture is asserted.
\subsection*{An exact compensated-weight test at the radial face}

The issue tested here is the following specific proposed inference: an
untwisted modulus partial-sum bound of size $X^\epsilon$, followed by
Abel summation for each dual vector and absolute summation over those
vectors, might retain only an $X^\epsilon$ loss because the oscillatory
integral compensates for its inverse-modulus phase.  The following
calculation shows precisely what that compensation does for one
family of vectors in the present problem.  It is a statement about
the full analytic weights, not a lower bound for the arithmetic error.

Write $e(x)=\exp(2\pi i x)$, and set
\[
 P=\{(s,t,z):0<z<1,\quad0<t<s-4z,\quad s<17z/4\},\qquad
 F(s,t,z)=s^2-3t^2-18z^2.
\]
Boundary conventions have no effect on the integrals below; the
original conventions may be retained when evaluating weights at
lattice points.  Fix a nonnegative $C^\infty$ function $\chi$ such that
\[
 0\leq\chi\leq1,\quad \chi(v)=1\ (|v|\leq1/128),\quad
 \chi(v)=0\ (|v|\geq1/64).
\]
For example, let $b(t)=\exp(-1/t)$ for $t>0$, $b(t)=0$ otherwise,
put $\theta(t)=b(t)/(b(t)+b(1-t))$, and take
$\chi(v)=1-\theta(128|v|-1)$.  This expression is smooth also at zero,
because it is constant in a neighborhood of zero.  Multiplying the
sharp weight by $\chi(F(y))$ preserves every counted solution for
$X\geq16$: on $F(Xy)=-2$ one has
$|F(y)|=2/X^2\leq1/128$.

We use exactly Heath-Brown's kernel, not a surrogate amplitude.
With
\[
 w_0(t)=\begin{cases}\exp(-(1-t^2)^{-1}),&|t|<1,\\0,&|t|\geq1,
 \end{cases}\quad
 c_0=\int_{-1}^1w_0(t)\,dt,\quad
 \omega(r)=4c_0^{-1}w_0(4r-3),
\]
the kernel is
\[
 h(r,v)=\sum_{j\geq1}\frac{\omega(rj)-\omega(|v|/(rj))}{rj}.
\]
These are the definitions in Heath-Brown (1996), Section 3; his
Theorem 2 supplies the Poisson formula for smooth compactly supported
weights, and Lemma 4 states the needed flatness in $v$.
The full primary author PDF was downloaded and its Theorems 1--2,
Section 3 and Lemma 4 read; see \cite{A5HeathBrown1996}.

Put $K=[2/3,3/4]$.  On the support of $\chi(F(y))$, for $X\geq16$,
\[
 |F(y)+2/X^2|\leq3/128<r/2\qquad(r\in K).
\]
Only $j=1$ can occur in the first sum defining $h$, and the second
sum vanishes.  Therefore the following identity is exact:
\begin{equation}\label{B:compensated-flat-h}
 h(r,F(y)+2/X^2)=\frac{\omega(r)}r\qquad
 (r\in K,\ \chi(F(y))\ne0).
\end{equation}
For $k\geq1$ define the normalized full integral and its Abel weight by
\[
 J_k(r)=\int_P\chi(F(y))h(r,F(y)+2/X^2)
                 e(-(s+kz)/r)\,dy,\qquad
 W_k(r)=r^{-1/2}J_k(r).
\]
By \eqref{B:compensated-flat-h}, these functions are independent of
$X$ throughout $K$.

\paragraph{Exact analytic conclusion.}
There is a fixed number $L>0$ such that, as $k\to\infty$,
\begin{equation}\label{B:compensated-norm}
 \|W_k\|_{L^\infty(K)}\asymp k^{-1},\qquad
 \operatorname{Var}_K W_k=L+O(k^{-1}).
\end{equation}
Consequently
\[
 \sum_{10\leq k\leq H}
 \bigl(\|W_k\|_\infty+\operatorname{Var}_K W_k\bigr)
 =LH+O(\log(2H)).
\]
The same linear lower bound holds for standard smooth approximations
to the sharp weight, for a fixed positive proportion of the frequencies
$k\leq H$; a precise statement is given below.

\paragraph{Proof, including the nonzero face coefficient.}
Write $s=zu$, $t=zv$.  The Jacobian is $z^2$, the fixed triangular
cross-section is
\[
 D=\{(u,v):4<u<17/4,\ 0<v<u-4\},\qquad
 f(u,v)=u^2-3v^2-18,
\]
and $F(zu,zv,z)=z^2f(u,v)$.  Thus
\[
 J_k(r)=\frac{\omega(r)}r
 \int_D\int_0^1 z^2\chi(z^2f(u,v))
               e(-(k+u)z/r)\,dz\,du\,dv.
\]
Twice integrating the inner integral by parts gives, uniformly on $K$,
\begin{equation}\label{B:compensated-leading}
 J_k(r)=-\frac{\omega(r)}{2\pi i k}A(r)e(-k/r)+O(k^{-2}),
 \qquad
 A(r)=\int_D\chi(f(u,v))e(-u/r)\,du\,dv.
\end{equation}
The lower endpoint contributes zero because $z^2\chi(z^2f)$ and its
first derivative vanish at zero.  The replacement
$(k+u)^{-1}=k^{-1}+O(k^{-2})$ is uniform on the bounded triangle.
All constants depend only on the displayed fixed cutoff.

The top-face coefficient cannot vanish by phase cancellation.  Indeed,
for $u_0=33/8$, one has $|u-u_0|\leq1/8$ on $D$, and hence
\[
 \Re\{e(u_0/r)A(r)\}
 \geq \cos(3\pi/8)\int_D\chi(f(u,v))\,du\,dv>0.
\]
For an explicit positive lower bound use the rational rectangle
\[
 R=[8489/2000,8491/2000]\times[159/2000,161/2000]\subset D.
\]
On this rectangle $|f|<1/128$, so $\chi(f)=1$, and
$\operatorname{area}(R)=10^{-6}$.  It follows that
\begin{equation}\label{B:compensated-A-positive}
 |A(r)|\geq10^{-6}\cos(3\pi/8)\qquad(r\in K).
\end{equation}

To control differentiation of the remainder without an unsupported
differentiation of big-$O$ notation, use Euler's identity.  Set
$a(r)=\omega(r)/r$ and
\[
 B_k(r)=\int_P F(y)\chi'(F(y))e(-(s+kz)/r)\,dy.
\]
Since
\[
 r\partial_r e(-(s+kz)/r)=-y\cdot\nabla_y e(-(s+kz)/r),\qquad
 \operatorname{div}(y\chi(F))=3\chi(F)+2F\chi'(F),
\]
the divergence theorem gives the exact relation
\begin{equation}\label{B:compensated-Euler}
 rJ_k'(r)=\left(\frac{ra'(r)}{a(r)}+3\right)J_k(r)
       -a(r)A(r)e(-k/r)+2a(r)B_k(r).
\end{equation}
The three angular faces satisfy $y\cdot\nu=0$ and have no flux.
The face $z=1$ has $y\cdot\nu=1$ and produces exactly the displayed
top-face term.  The vertex at zero has zero flux; alternatively apply
the theorem first to $z\geq\delta$ and let $\delta\downarrow0$.
One radial integration by parts gives $J_k=O(k^{-1})$ and
$B_k=O(k^{-1})$, uniformly in $r$.  Hence
\[
 J_k'(r)=-\frac{\omega(r)}{r^2}A(r)e(-k/r)+O(k^{-1}),\qquad
 W_k'(r)=-\frac{\omega(r)}{r^{5/2}}A(r)e(-k/r)+O(k^{-1}).
\]
Thus
\begin{equation}\label{B:compensated-L}
 L=\int_{2/3}^{3/4}\frac{\omega(r)}{r^{5/2}}|A(r)|\,dr>0,
\end{equation}
which proves \eqref{B:compensated-norm}.  The same radial integration
by parts, after differentiating twice under the integral, gives
$\|W_k''\|_\infty=O(k)$: the differentiated amplitude and its first
radial derivative are $O(k^2)$, while the phase derivative is
$\asymp k$.  This observation also verifies the discrete assertion
below.

The positivity can be certified without numerical quadrature of
$A$.  Since $c_0\leq2e^{-1}$ and
$w_0(4r-3)\geq e^{-9/8}$ on $K$, one has
$\omega(r)\geq2e^{-1/8}$.  Therefore
\begin{equation}\label{B:compensated-explicit-lower}
 L\geq\frac1{12}\,2e^{-1/8}\left(\frac43\right)^{5/2}
           10^{-6}\cos(3\pi/8)>0.
\end{equation}
The rational rectangle checks, Euler identities and a rigorous Arb
enclosure of the last positive number are recorded by
\path{code/B_A5_compensated_weight.py} and its stdout.

\paragraph{The actual discrete modulus variation.}
Let $\operatorname{Var}_{X,K} W$ be the sum of
$|W((q+1)/X)-W(q/X)|$ over consecutive integer mesh points lying in
$K$.  The elementary comparison of a differentiable function with its
piecewise linear interpolant gives
\[
 \left|\operatorname{Var}_{X,K}W_k-
                  \operatorname{Var}_K W_k\right|
 \ll X^{-1}(\|W_k'\|_\infty+\|W_k''\|_{L^1(K)})
 \ll (1+k)/X.
\]
For example, on each mesh interval replace $W_k'(r)$ by its left
endpoint value, and bound both resulting errors by the integral of
$|W_k''|$ times the mesh width.  The two incomplete endpoint intervals
cost $O(X^{-1}\|W_k'\|_\infty)$.
It follows that
\[
 \sum_{10\leq k\leq H}\operatorname{Var}_{X,K}W_k
       =LH+O(\log(2H)+H^2/X).
\]
In particular this is a loss of order $H$ when $H\asymp\sqrt X$.
Restricting $q$ to any one fixed residue class only changes the mesh
width by a fixed factor and gives the same conclusion.

\paragraph{Stability under boundary smoothing.}
Let $w=\mathbf1_P\chi(F)$, and let $w_H$ be any smooth approximation
supported in a fixed compact set within $|F|\leq1/64$, with
\[
 \|w_H-w\|_{L^1(\R^3)}\leq C/H.
\]
Such approximations are obtained by mollifying $\mathbf1_P$ at scale
$H^{-1}$ and then multiplying by $\chi(F)$; the polytope boundary has
finite area.  Define $J_{H,k}$ and $W_{H,k}$ by replacing $w$ with
$w_H$ in the full integrals.  Identity
\eqref{B:compensated-flat-h} still holds exactly.  Differentiating its
finite integral gives
\[
 \|W_{H,k}-W_k\|_\infty\ll H^{-1},\qquad
 \|W_{H,k}'-W_k'\|_\infty\ll(1+k)/H.
\]
Consequently there are fixed $\eta,c>0$ and $k_0$ such that, for
$k_0\leq k\leq\eta H$, $H=o(X)$, and $X$ sufficiently large,
\[
 \operatorname{Var}_{X,K}W_{H,k}\geq c.
\]
Indeed compare discrete variations by the integral of
$|W_{H,k}'-W_k'|$, choose $k_0$ so that the $O(k^{-1})$ error is
less than $L/4$, then choose $\eta$ so that the smoothing error is
less than $L/4$, and finally use $H/X\to0$.  The upper bound is
also fixed on this frequency range.  Thus the sum of these full
smoothed variation norms is $\asymp H$.

\paragraph{What this proves and what it leaves open.}
For $c_k=(1,0,k)$ the good-modulus dual form in the exact Gauss
calculation is
\[
 H(c_k)=k^2-18.
\]
For $k\geq10$ it lies strictly between $(k-1)^2$ and $k^2$ and is
not a rational square.  In the rational normalization of the Sali\'e
parameters the discriminant is $4mn=H(c_k)/9$, also not a square.
This avoids the evident square-discriminant family $c=(0,0,k)$.
We do not use this observation as a substitute for identifying the
precise half-integral multiplier and its complete unary-theta
projection; that arithmetic step remains separate.

On a modulus block $q=Xr$, the exact Poisson summand has the scale
\[
 X^{-2}q^{-3}S_q(c)I_q(c),\qquad I_q(c_k)=X^3J_k(r).
\]
After the three Gauss factors, writing the arithmetic sequence in the
normalization $T_q(c_k)/q$ leaves the analytic weight
\[
 q^{-1/2}I_q(c_k)=X^{5/2}W_k(q/X).
\]
Grant, even before completing the multiplier identification, an
$O_\epsilon((Xk)^\epsilon)$ bound for every required partial sum of
the theta-subtracted arithmetic sequence on this block.  Abel
summation and absolute summation over $k$ then incur the weight norm
just calculated: their direct bound has scale
$X^{1/2+\epsilon}H$, not $X^{1/2+\epsilon}$, on this family alone.
This statement concerns the bound delivered by that procedure; it is
not a lower bound for the true arithmetic sum.  In the common
balancing heuristic with a boundary cost $X/H$, it predicts the
exponent $3/4$, not the requested $1/2$; no full $3/4$ theorem is
claimed here because the other dual frequencies and local
arithmetic factors have not been bounded by this calculation.

The missing improvement could come from cancellation between dual
frequencies, an arithmetic transformation retaining their phases,
or another organization of the count.  This calculation does not
rule out such arguments, nor does it prove that Sun's conjecture
cannot imply the desired result.  It settles only the proposed free
compensation in the termwise Abel-summation argument, while using
the actual radial face and the full Heath-Brown kernel.
